\SourceLeaf{242}%
tas \(x,y,z\) tamquam functiones indeterminatarum \(p,q\) exhibeant, sed sufficere expressionem generalem pro magnitudine cuiusvis elementi linearis. Progrediamur ad aliquot applicationes huius gravissimi theorematis.

Supponamus, superficiem nostram curvam explicari posse in aliam superficiem, curvam seu planam, ita ut cuivis puncto prioris superficiei per coordinatas \(x,y,z\) determinato respondeat punctum determinatum superficiei posterioris, cuius coordinatae sint \(x',y',z'\). Manifesto itaque \(x',y',z'\) quoque considerari possunt tamquam functiones indeterminatarum \(p,q\), unde pro elemento \(\BareRoot{(\mathrm{d}x'^2+\mathrm{d}y'^2+\mathrm{d}z'^2)}\) prodibit expressio talis
\[
\BareRoot{(E'\mathrm{d}p^2+2F'\mathrm{d}p\mathord{.}\mathrm{d}q+G'\mathrm{d}q^2)}
\]
denotantibus etiam \(E',F',G'\) functiones ipsarum \(p,q\). At per ipsam notionem \textit{explicationis} superficiei in superficiem patet, elementa in utraque superficie correspondentia necessario aequalia esse, adeoque identice fieri
\[
E=E',\quad F=F',\quad G=G'
\]
Formula itaque art. praec. sponte perducit ad egregium

\textsc{Theorema.} \textit{Si superficies curva in quamcunque aliam superficiem explicatur. mensura curvaturae in singulis punctis invariata manet.}

Manifesto quoque \textit{quaevis pars finita superficiei curvae post explicationem in aliam superficiem eandem. curvaturam integram retinebit.}

Casum specialem, ad quem geometrae hactenus investigationes suas restrinxerunt, sistunt superficies in planum explicabiles. Theoria nostra sponte docet, talium superficierum mensuram curvaturae in quovis puncto fieri \(=0\), quocirca, si earum indoles secundum modum tertium exprimitur, ubique erit
\[
\frac{\mathrm{d}\mathrm{d}z}{\mathrm{d}x^2}\mathord{.}\frac{\mathrm{d}\mathrm{d}z}{\mathrm{d}y^2}-(\frac{\mathrm{d}\mathrm{d}z}{\mathrm{d}x\mathord{.}\mathrm{d}y})^2=0
\]
quod criterium, dudum quidem notum, plerumque nostro saltem iudicio haud eo rigore qui desiderari posset demonstratur.
\Article{13}
Quae in art. praec. exposuimus, cohaerent cum modo peculiari superficies considerandi, summopere digno, qui a geometris diligenter excolatur. Scilicet quatenus superficies consideratur non tamquam limes solidi, sed tamquam soli%
